% compendium/volatility-trading-premia/vix-variance-dividends.tex
% Entry migrated on 2026-10-16 from the legacy macro chapter (slice 37, 02_transmission_maps.tex), owned by
% vix-correlation-dividends. Gliner (2014) is cited by chapter, section and PDF page of the ebook.
% Dispositions in notes/migration-slices.md, slice 37.

\section{VIX, variance swaps and dividend swaps}

\begin{qfdefinition}[The VIX and its futures curve, variance swaps (strike, vega notional, variance units and payoff) and dividend swaps, with Gliner's pricing statement corrected]\label{def:vix-variance-dividend-swaps}
\qffield{Statement} \emph{VIX.} Gliner describes the VIX as derived from S\&P 500 options ``one month out'', negatively correlated with equity returns, able to move ``from the mid-teens to 30 in a matter of days'' in sell-offs; VIX futures ``typically trade in contango'' and move to backwardation when spot spikes; VIX options are in effect options on the volatility of volatility; and other indices include V2X, VXN, VXD, RVX, GVZ and OVX \citep[chap.~8, sec.\ ``The Volatility Index (VIX)'', PDF pp.~222--225]{s7-gliner}. \emph{Variance swaps} give ``pure variance exposure'' without delta hedging and without the interest-rate and dividend exposure of options; they are quoted by strike, vega notional, variance units and maturity (pp.~225--226). Variance units $=$ vega notional$/(2K)$ and payoff $=(\sigma_R^2-K^2)\times$ variance units: with $K=24$, vega notional \$1,000,000 and realised volatility 29, there are 20,833 units and a payoff of \$5,520,833 (p.~226). \emph{Dividend swaps} let one take a view that an index's dividends will be higher or lower than a fixed amount, OTC for the S\&P 500 and listed for the Euro Stoxx 50 (sec.\ ``Dividend Swaps'', p.~227).

\qffield{Notation} $K$ strike in volatility points; $\sigma_R$ realised volatility over the life, annualised from daily log returns, in points; vega notional is the approximate P\&L for a one-point move in realised volatility near the strike.

\qffield{Conditions} \emph{The strike is fixed at inception from option prices.} Gliner writes that variance swaps ``are priced with realized volatility, which gives the buyer slightly more edge'' (p.~225); the payoff depends on realised variance, but the strike is set at the start from the market's implied variance and is usually above subsequent realised variance (the variance risk premium), so the long typically pays: with strike 24 and realised 20 the long loses \$3,666,667 (\nolinkurl{notes/artefacts/multi-asset-systematic-macro/slice37-checks.py}, section~13). The payoff is convex in volatility, linear in variance.

\qffield{Practice} \emph{Daily moves.} With $\sqrt{252}=15.87$, a volatility of 16 corresponds to a daily standard deviation of about 1 percent and 24 to about 1.5 percent (Gliner, p.~226; checks, section~13); that is a standard deviation, not the ``expected daily move'', whose mean absolute size under normality is about 0.8 of it. \emph{Dividend swap example} (p.~227): with the 2014 S\&P 500 dividend quoted at \$31 against an index of 1,350 (an implied yield of 2.3 percent), a buyer of 500,000 units gains $(36-31)\times500{,}000=\$2.5$ million if the market moves to \$36; Gliner's ``baskets of 100,000'' sits uneasily with that arithmetic, which counts one index unit per basket.

\qffield{Pitfalls} The VIX is computed from a strip of out-of-the-money S\&P 500 option prices as a 30-day variance-swap rate, not by extrapolating implied volatilities as Gliner puts it; Cboe's methodology \citep[catalogue entry only, not read]{gap-cboe-vix6} is the reference (backlog). Gliner's ``V2X Index---Euro Stoxx 500'' should read Euro Stoxx 50. ``VIX options provide asymmetric hedges'' ignores that VIX options are priced off VIX futures, not spot.

\qffield{See also} \hyperref[def:equity-indices-derivatives]{equity indices}; \hyperref[meth:equity-macro-leading-indicators]{AUD volatility and the VIX}; \hyperref[def:volatility-components]{volatility components}.

\qffield{Further reading} Variance swaps \citep{wide-demeterfi-varswaps}; the VIX premium \citep{gap-vix-premium}; dividend strips \citep{gap-dividend-strips}.

\qffield{Index} VIX; volatility index; VIX futures; contango; backwardation; VIX options; volatility of volatility; V2X; VXN; OVX; variance swap; variance strike; vega notional; variance units; realised variance; variance risk premium; dividend swap; implied dividend

\qffield{Status} Verification: partly checked 2026-10-16 (Gliner's text, chap.~8, PDF pp.~222--227; the variance-swap and daily-move arithmetic recomputed in the slice's check script; Cboe's methodology not read). Stability: draft. Currency: Gliner's examples as of 2013.
\end{qfdefinition}
